\section{Introduction}
\label{sec:intro}

\IEEEPARstart{A}{} microservice application is assembled from many small services that talk to
each other over the network and can be built, deployed and scaled one at a time
\cite{dragoni2017,jamshidi2018}. Teams adopt the style for that independence, and pay for it
in operations. A systematic review of grey literature records the pains practitioners report
alongside the gains \cite{soldani2018}, and performance engineering for these systems has its
own list of open research problems \cite{heinrich2017}. Failure is near the top of that list.
A study of hundreds of cloud service outages draws lessons about why such systems stop
working \cite{gunawi2016}, and an industrial survey of microservice systems reports that
faults are slow to diagnose \cite{zhou2021}.

Operators would like to know in advance which services deserve hardening first. Injecting
faults into every service answers that directly, but each answer costs a run per service and
per fault type. A cheaper answer can be read from the dependency graph. Graph measures have
long served as indicators of importance \cite{freeman1978}, and network measures computed on
software dependency graphs have predicted which components accumulate defects
\cite{zimmermann2008}. Applied to a service graph, the familiar reading is that a service
calling many others is a hub, and a hub is dangerous. This project started from that reading.
Its original metric, hybrid criticality, scales out-degree centrality by fan-out, the number
of services a service calls directly, on the argument that a service with many callees cuts
many links when it fails.

There is a simple argument in the other direction. When a service fails, the services left
waiting are the ones that call it; its callees simply receive less traffic. On that argument
the number of services that degrade should follow a service's transitive callers, its
ancestors in the graph, and a fan-out ranking should go most wrong on exactly the services it
rates highest. The composition service of the Social Network benchmark shows the clash in one
place. It calls \val{composepost_out_degree} services directly, and the fan-out metric ranks
it first %[id:composepost_rank_hybrid]
of the \val{composepost_set_size} services analysed. It has a single caller,
%[id:composepost_ancestor_count]
so a ranking by callers puts it last. %[id:composepost_rank_ancestor_count]

This paper tests the caller-based prediction against measurement, with the fan-out metric
kept as a comparator, in a preregistered \cite{nosek2018} fault-injection campaign on two
applications of the DeathStarBench suite \cite{gan2019}. The caller-based hypothesis was
formed after a pilot study and locked, together with every analysis choice, before any
campaign data existed. All \val{campaign_runs_total} planned runs completed. Each analysable
service received a process kill and an injected network delay, \val{campaign_reps_per_cell}
times each. Dependency graphs were rebuilt from traces, latency was measured on the spans
each service emits itself, and every degraded service was classified by its graph relation
to the faulted one, so that degraded callers and degraded callees were counted separately.

The preregistered hypothesis was not confirmed. On Social Network under kill faults, the
caller count and the number of degraded callers correlated at $\rho = \val{primary_rho}$, but
none of the \val{conf_family_size} tests in the confirmatory family survived correction for
multiple testing, and a replication under latency faults gave the same answer.
%[id:conf_family_significant] %[id:repl_family_significant]
A large part of that correlation is definitional, because the outcome counts a subset of the
predictor and most services degraded every caller they had. The fan-out metric fared no
better: its correlation with degraded callers was negative for both applications and both
fault types, %[id:h1b_all_negative]
and in no case distinguishable from zero. %[id:h1b_any_distinguishable]
The plainest result needs no test. Degradation reached callers and almost never callees: no
descendant of the faulted service degraded in \val{descendant_affected_runs_zero} of
\val{campaign_runs_total} runs.

The paper makes four contributions:
\begin{itemize}
\item A preregistered test of whether caller count, or a fan-out-weighted centrality metric,
tracks measured failure impact in two benchmark applications, reported with its null result
and with the ceiling that the design placed on any correlation.
\item Evidence from \val{campaign_runs_total} runs under two fault types that degradation
travels to the callers of a failed service and almost never to its callees.
\item A measurement pitfall in trace-based latency attribution, in which a request's
end-to-end latency is charged to every service on its path, together with its effect in a
pilot and the per-service attribution that removes it.
\item An artifact in which a script regenerates every number in the paper from the logged
campaign data.
\end{itemize}

Section~\ref{sec:related} reviews related work and Section~\ref{sec:problem} states the problem
and the hypotheses. Section~\ref{sec:method} describes the testbed and the method,
Section~\ref{sec:results} reports the results and Section~\ref{sec:discussion} interprets them.
Section~\ref{sec:threats} discusses threats to validity and Section~\ref{sec:conclusion}
concludes.
